% ============================================================================
\clearpage
\section{Protocolo de evidencia}\label{sec:evidencia}
% ============================================================================

\ficha{Registro previo de objetivos, transformaciones y variantes; series de pronósticos y decisiones
fuera de muestra.}
{Particiones temporales con purga, comparación incremental, pruebas de pronóstico y de política,
corrección por selección y estrés.}
{Estimaciones de mejora con intervalos por bloques, veredicto por objetivo y horizonte, informe de
estrés.}
{Umbrales fijados antes del holdout; ``evidencia insuficiente'' cuando el intervalo no se separa de cero.}

\subsection{Registro previo}\label{sec:preregistro}

Antes de probar nada se registran los objetivos, las transformaciones de variables y el número de variantes
que se evaluarán. Cada experimento posterior queda en el registro, haya funcionado o no: el número total
de pruebas es un insumo de la corrección por selección (\secref{sec:seleccion}).

\subsection{Particiones temporales, purga y tamaño efectivo}\label{sec:particiones}

Se separan cuatro tramos con funciones distintas: entrenamiento, ajuste de hiperparámetros, calibración y
prueba final. La validación es \emph{walk-forward}, sin particiones aleatorias, y con purga de las
observaciones cuya etiqueta cruza una frontera entre tramos \citep{lopezdeprado2018}. Cuando el esquema
pueda reintroducir observaciones correlacionadas se añade un embargo (\figref{fig:walkforward}).

\begin{figure}[H]
\centering
\resizebox{\linewidth}{!}{% Walk-forward con purga y embargo; tramos de entrenamiento, ajuste, calibración y prueba.
\begin{tikzpicture}[x=1cm,y=1cm,
  tramo/.style={minimum height=5.2mm,inner sep=0pt,font=\sffamily\tiny,text=tinta,anchor=west},
  purga/.style={minimum height=5.2mm,inner sep=0pt,anchor=west,pattern=north east lines,pattern color=tenue,
    draw=tenue,line width=0.3pt}]
\def\bloque#1#2{% #1 = fila (y), #2 = longitud del entrenamiento
  \node[tramo,fill=qazul!40,minimum width=#2 cm] (e) at (0,#1) {entrenamiento};
  \node[purga,minimum width=3.5mm] (g1) at (e.east) {};
  \node[tramo,fill=qnaranja!40,minimum width=1.35cm] (h) at (g1.east) {\ ajuste hiperp.};
  \node[purga,minimum width=3.5mm] (g2) at (h.east) {};
  \node[tramo,fill=qaqua!40,minimum width=1.35cm] (c) at (g2.east) {\ calibración};
  \node[purga,minimum width=3.5mm] (g3) at (c.east) {};
  \node[tramo,fill=tinta2!45,minimum width=1.2cm,text=white] (p) at (g3.east) {\ prueba};}
\node[titulo,anchor=east] at (-0.15,0) {Ventana 1};
\bloque{0}{3.2}
\node[titulo,anchor=east] at (-0.15,-0.75) {Ventana 2};
\bloque{-0.75}{4.4}
\node[titulo,anchor=east] at (-0.15,-1.5) {Ventana 3};
\bloque{-1.5}{5.6}
% Prueba final reservada
\fill[pattern=crosshatch,pattern color=eje] (11.35,0.35) rectangle (14.3,-1.85);
\draw[tinta2,line width=0.6pt] (11.35,0.35) rectangle (14.3,-1.85);
\node[nota,text=tinta,align=center,fill=white,inner sep=2pt] at (12.82,-0.75) {\textbf{prueba final}\\ se abre una\\ sola vez};
% Eje y leyenda
\draw[flecha,draw=tinta2] (0,-2.2) -- (14.5,-2.2);
\node[nota,anchor=north east] at (14.5,-2.25) {tiempo};
\node[purga,minimum width=3.5mm] at (0,-2.75) {};
\node[nota,anchor=west] at (0.45,-2.75) {purga y embargo: se eliminan observaciones cuya etiqueta cruza la frontera, más un margen};
\end{tikzpicture}
}
\caption{Validación walk-forward. Cada ventana entrena con el pasado, ajusta, calibra y prueba en tramos
sucesivos separados por purga y embargo; la prueba final queda reservada hasta el final.}
\label{fig:walkforward}
\end{figure}

La dependencia reduce la información disponible más de lo que sugiere el número de filas. Con etiquetas
solapadas de $h$ sesiones, la autocorrelación inducida es $1-k/h$ y el tamaño efectivo es
aproximadamente $n/h$ (\figref{fig:obs_efectivas}); con $N$ activos de correlación media $\bar\rho$, el
número de activos independientes equivalentes es $N/[1+(N-1)\bar\rho]$ (\figref{fig:activos_efectivos}).
Por eso los intervalos se calculan con bloques temporales y teniendo en cuenta la dependencia entre activos.

\grafica[0.54]{s7_obs_efectivas}{Observaciones efectivas en 500 sesiones según el horizonte de la
etiqueta: un pronóstico a 20 sesiones tiene del orden de 25 observaciones independientes.}{fig:obs_efectivas}

\grafica[0.54]{s7_activos_efectivos}{Activos independientes equivalentes según la correlación media. Cinco
acciones con correlación 0.6 equivalen a menos de dos.}{fig:activos_efectivos}

\subsection{Comparación incremental}\label{sec:incremental}

La pregunta de investigación se responde con una escalera de conjuntos de información
(\figref{fig:incremental}): A, precio, volumen y volatilidad; B, A más la distribución $\Q$ actual; C, B
más cambios simples; D, C más transporte e innovaciones; E, D más aceleración y regímenes. Se compara
dentro de una misma familia de predicción para atribuir el cambio a la información adicional y no al
aumento simultáneo de capacidad del modelo.

\begin{figure}[H]
\centering
\resizebox{\linewidth}{!}{% Comparación incremental A-E: cada paso añade información, no capacidad del modelo.
\begin{tikzpicture}[x=1cm,y=1cm,
  paso/.style={draw=qazul,fill=qazul!#1,rounded corners=2pt,line width=0.6pt,anchor=south west,
    text width=2.55cm,align=center,font=\sffamily\scriptsize,text=tinta,inner sep=3pt}]
\node[paso=6,minimum height=1.0cm] (A) at (0,0) {\textbf{A}\\ precio, volumen\\ y volatilidad};
\node[paso=12,minimum height=1.55cm] (B) at (2.95,0) {\textbf{B = A +}\\ distribución $\Q$ actual};
\node[paso=18,minimum height=2.1cm] (C) at (5.9,0) {\textbf{C = B +}\\ cambios simples};
\node[paso=24,minimum height=2.65cm] (D) at (8.85,0) {\textbf{D = C +}\\ transporte e\\ innovaciones};
\node[paso=30,minimum height=3.2cm] (E) at (11.8,0) {\textbf{E = D +}\\ aceleración y\\ regímenes};
\foreach \a/\b in {A/B,B/C,C/D,D/E} \draw[flecha] (\a.east) ++(0,0.3) -- ++(0.36,0);
\draw[guia,line width=0.7pt] (-0.1,-0.05) -- (14.5,-0.05);
\node[nota,anchor=north] at (7.2,-0.15) {misma familia de predicción e hiperparámetros equivalentes en cada comparación:
  la mejora se atribuye a la información añadida, no a más capacidad};
\end{tikzpicture}
}
\caption{Comparación incremental A--E. La hipótesis principal corresponde al paso de C a D.}
\label{fig:incremental}
\end{figure}

\subsection{Dos pruebas independientes: pronóstico y política}\label{sec:dos_pruebas}

Una señal puede predecir mejor y no pagar sus costos. Por eso hay dos pruebas separadas:

\begin{itemize}
\item \textbf{Mejora de pronóstico}: diferencial de reglas de puntuación \eqref{eq:diferencial} por
objetivo y horizonte.
\item \textbf{Mejora de política de posición}: rendimiento neto, drawdown, pérdidas de cola, rotación y
exceso frente a un benchmark de riesgo comparable, contra las referencias de la \tabref{tab:referencias}.
\end{itemize}

\begin{table}[H]
\centering\small
\caption{Referencias para la prueba de política. Se alinean moneda, aportaciones, horarios y exposición.}
\label{tab:referencias}
\begin{tabularx}{\linewidth}{@{}l Y@{}}
\toprule
\encab{Referencia} & \encab{Qué controla}\\
\midrule
Mantener las posiciones & El valor de no hacer nada, incluidos costos evitados.\\
SPY con dividendos & La exposición pasiva al mercado.\\
Control de volatilidad sencillo & Si la mejora es solo gestión de riesgo genérica.\\
Pesos iguales & Para selección de acciones: si el modelo supera una asignación ingenua.\\
\bottomrule
\end{tabularx}
\end{table}

\subsection{Sesgo de selección}\label{sec:seleccion}

Probar muchas variantes y quedarse con la mejor infla el resultado aunque ninguna tenga ventaja real. Bajo
la hipótesis nula, el máximo esperado de $N$ ratios de Sharpe independientes con varianza $V$ es
aproximadamente \citep{bailey2014}
\begin{equation}
\widehat{SR}_0=\sqrt{V}\Bigl[(1-\gamma)\,\Phi^{-1}\!\Bigl(1-\tfrac1N\Bigr)+\gamma\,\Phi^{-1}\!\Bigl(1-\tfrac{1}{Ne}\Bigr)\Bigr],
\label{eq:sr0}
\end{equation}
con $\gamma$ la constante de Euler--Mascheroni. El Sharpe deflactado es la probabilidad de que el Sharpe
verdadero supere ese umbral, corrigiendo por asimetría $\hat\gamma_3$ y curtosis $\hat\gamma_4$ de $T$
rendimientos:
\begin{equation}
\mathrm{DSR}=\Phi\!\left(\frac{(\widehat{SR}-\widehat{SR}_0)\sqrt{T-1}}{\sqrt{1-\hat\gamma_3\,\widehat{SR}+\frac{\hat\gamma_4-1}{4}\,\widehat{SR}^2}}\right).
\label{eq:dsr}
\end{equation}
Se usa como diagnóstico adicional, junto con la evaluación fuera de muestra y el bootstrap por bloques; no
como certificado definitivo (figuras~\ref{fig:maximo_sharpe} y~\ref{fig:sharpe_deflactado}).

\grafica[0.56]{s7_maximo_sharpe}{Sharpe anual máximo esperado entre $N$ estrategias sin ninguna ventaja
real. Con dos años de datos diarios y cien variantes, el mejor resultado ronda 1.8 por puro azar; los
puntos confirman la aproximación por simulación.}{fig:maximo_sharpe}

\grafica[0.56]{s7_sharpe_deflactado}{Sharpe deflactado de un mismo resultado observado según el número
de variantes probadas (supuestos en el eje).}{fig:sharpe_deflactado}

\subsection{Intervalos por bloques y evidencia insuficiente}\label{sec:bootstrap}

La mejora media del diferencial de pérdida se estima con bootstrap de bloques móviles \citep{kunsch1989},
que conserva la dependencia dentro de cada bloque. Si el intervalo no permite distinguir la mejora de cero,
el veredicto es \emph{evidencia insuficiente}, no ``casi significativo''
(figuras~\ref{fig:bootstrap_distinguible} y~\ref{fig:bootstrap_insuficiente}).

\grafica[0.56]{s7_bootstrap_distinguible}{Distribución bootstrap de la mejora media del modelo D sobre el C
(serie sintética con autocorrelación). El intervalo del 90\% excluye el cero.}{fig:bootstrap_distinguible}

\grafica[0.56]{s7_bootstrap_insuficiente}{Una mejora media menor con la misma variabilidad: el intervalo
contiene el cero y el veredicto es evidencia insuficiente.}{fig:bootstrap_insuficiente}

\subsection{Pruebas de estrés}\label{sec:estres}

\begin{table}[H]
\centering\small
\caption{Pruebas de estrés. Ningún horario ni umbral se selecciona con el periodo final.}
\label{tab:estres}
\begin{tabularx}{\linewidth}{@{}l Y@{}}
\toprule
\encab{Perturbación} & \encab{Qué se observa}\\
\midrule
Choques de mercado y de correlación & CVaR, límites y cambio de recomendación.\\
Spreads más amplios & Ancho de bandas, estados de confianza y abstención.\\
Retrasos y cotizaciones incompletas & Que la decisión use solo datos disponibles en el corte.\\
Interrupciones y degradación del proveedor & Que un fallo no se convierta en orden.\\
Resultados trimestrales y cambios de dividendo & Conversión americana, forward y deformación por plazo.\\
Revisiones de datos & Reproducibilidad con los datos conocidos en cada momento.\\
Órdenes parcialmente ejecutadas & Reconciliación de posiciones y caja.\\
\bottomrule
\end{tabularx}
\end{table}

\subsection{Regla de avance}\label{sec:avance}

\begin{criterio}
Un modelo avanza si mejora la métrica definida de antemano, mantiene una calibración aceptable y produce
decisiones estables frente a perturbaciones plausibles. Los umbrales numéricos se fijan antes del holdout
con el presupuesto de riesgo acordado.
\end{criterio}

\begin{noconcluir}
Llegar a diciembre con el sistema operativo no equivale a haber demostrado ventaja estadística. Un sistema
que funciona sin fallas y una hipótesis confirmada son dos entregables distintos.
\end{noconcluir}
